\section{Sampling contrasts for heuristics and biases}
\label{sec:annex-p19-rationality-2-4}

\begin{annexitem}{annex:p19.r24.exposure}{thm:p19-r24-exposure}{Selection weights and the nonidentification construction}
\begin{proof}
Let \(I\) denote admission to the exposure record. The stipulated
conditional selection probabilities give
\(P(\text{category }i,I)=p_i s_i\) and
\(P(I)=\sum_j p_j s_j>0\). Dividing yields the displayed
conditional proportions. The numerical example has selected masses
\(.009,.1\), total \(.109\), hence the two stated fractions
and reversed ordering. A sample estimate unbiased for the admitted
record's law has mean \(q_i\); substituting that record law for the
target law is not justified by unbiasedness for the former.

For the general construction, the set of indices with \(q_i>0\)
is nonempty. Choose
\[
 \lambda=\min\left(1,\min_{i:q_i>0}\frac{p_i}{q_i}\right)>0,
 \qquad s_i=\frac{\lambda q_i}{p_i}.
\]
Strict positivity of every \(p_i\) makes these definitions valid.
For a zero \(q_i\), \(s_i=0\); otherwise the choice of
\(\lambda\) ensures \(s_i\leq1\). Thus all weights are
admissible and \(\sum_i p_i s_i=\lambda\sum_i q_i=\lambda>0\).
The output proportion is \(\lambda q_i/\lambda=q_i\).
Consequently any observed profile on this finite domain is compatible
with appropriately selected exposure followed by unbiased recording.
An observationally identical profile can also be produced by an
unselected exposure law equal to \(q\). More specific claims about
the actual exposure, encoding or retrieval mechanism require additional
evidence; this construction merely establishes nonidentification from
the final profile alone.
\end{proof}
\end{annexitem}

\begin{annexitem}{annex:p19.r24.recognition-mixture}{thm:p19-r24-recognition-mixture}{The finite mixture and its comparison boundary}
\begin{proof}
There are \(N(N-1)/2\) unordered pairs: \(n(N-n)\)
mixed pairs, \(n(n-1)/2\) both-recognized pairs, and
\((N-n)(N-n-1)/2\) both-unrecognized pairs. Their conditional
accuracies are respectively \(\alpha,\beta,1/2\). Weighting
by the pair counts yields the stated formula; empty classes
have zero weight regardless of their assigned parameter versions.
At \(n=N\), the value is \(\beta\). At \(n=N-1\),
it is \(((N-2)\beta+2\alpha)/N\), proving the difference
formula. Thus \(\alpha>\beta\) suffices for a smaller
recognition amount to beat complete recognition. Conversely,
if \(\alpha\leq\beta\) and \(1/2\leq\beta\), every
conditional accuracy in the mixture is at most \(\beta\),
so no recognition amount can beat it. Both directions use
the fixed-validity assumption for comparisons across amounts.
For \(N=100,n=50\), the numerator is
\(2(50)(50)(.8)+(50)(49)(.6)+(50)(49)/2=6695\),
and the denominator is 9900. At complete recognition the
accuracy is .6. If \(\beta<1/2\), even the no-recognition
guessing baseline beats complete recognition, which is why
the additional boundary is part of the equivalence statement.
\end{proof}
\end{annexitem}

\begin{annexitem}{annex:p19.r24.priority-invariance}{thm:p19-r24-priority-invariance}{An excluded influence at a stopped first stage}
\begin{proof}
The first-stage stopping predicate and the choice made at that
stage depend only on the fixed outcomes and declared threshold.
The stipulated probability changes leave that predicate true,
leave its preferred minimum outcome unchanged, and leave the
pair in the applicability domain. Strict stopping therefore
prevents examination of a later probability cue and returns
the same choice in every such case. A systematic contrary
influence cannot be an output of this specified deterministic
procedure. Another procedure, response noise or a changed
interpretation may produce it; the invariance proof by itself
does not choose among those explanations. The argument
does not interchange rounded and unrounded thresholds.
\end{proof}
\end{annexitem}

\begin{annexitem}{annex:p19.r24.strings}{thm:p19-r24-strings}{Exact strings, fair marginals and observed samples}
\begin{proof}
Independence makes the probability of an exact string the product of
its six position probabilities. GBGBBG contains three B's and three
G's; BGBBBB contains five B's and one G. The stated products follow.
The coin strings contain three and one H respectively, yielding their
products. At fairness, each factor is \(1/2\), so any exact six-
position string has probability \(2^{-6}\), irrespective of its
apparent regularity. This is not the probability of every string with
three B's, which is the sum over twenty disjoint exact strings.

To refute inference from fair marginals alone, assign probability
\(1/2\) to GBGBBG and \(1/2\) to its positionwise complement
BGBGGB. Each position then has marginal probability \(1/2\) for
B, but the first target string has probability \(1/2\) and
BGBBBB has probability zero. The construction with HTHTTH and its
complement THTHHT gives the analogous coin counterexample. The
coordinates in these examples are dependent.

Finally, with independent fair sampling of one family, the counts
of the two birth strings can be \((1,0)\) with positive probability
\(1/64\), although their underlying probabilities are equal. Thus
equal probabilities do not require equal observed counts. The next
annex result supplies the conditioned larger-sample calculation.
\end{proof}
\end{annexitem}

\begin{annexitem}{annex:p19.r24.counts}{thm:p19-r24-counts}{The multinomial conditional distribution}
\begin{proof}
Let the remaining category have probability \(1-a-b\).
Independent sampling of \(N\) cases gives, for \(x+y\leq N\),
\[
 P(X=x,Y=y)=\frac{N!}{x!y!(N-x-y)!}
       a^x b^y(1-a-b)^{N-x-y}.
\]
The marginal count has
\(P(X=x)=\binom Nx a^x(1-a)^{N-x}\).
For an attainable \(x\), this denominator is positive. Dividing
and simplifying gives
\[
 P(Y=y\mid X=x)=\binom{N-x}{y}
 \left(\frac{b}{1-a}\right)^y
 \left(1-\frac{b}{1-a}\right)^{N-x-y},
       \quad 0\leq y\leq N-x.
\]
The parameter is in \([0,1]\) because \(a+b\leq1\), and
the denominator \(1-a\) is positive. Empty-product conventions
make the formula valid at \(a=0\), \(b=0\) and
\(a+b=1\) whenever conditioning is attainable. When \(N=x\),
the count has no remaining trials and is identically zero.
The binomial mean is \((N-x)b/(1-a)\), obtained by summing its
\(N-x\) Bernoulli indicators or differentiating the finite binomial
expansion. For the fair distinct strings, \(a=b=1/64\), so this
mean is \((N-x)/63\). Set \(x=72\) and solve
\((N-72)/63=72\) to obtain \(N=4608\). Since the conditioned
mean varies with \(N\), the stated first count alone cannot identify
it. Equal unconditioned expectations \(N/64\) for the two counts
are a different assertion and are not contradicted by this calculation.
\end{proof}
\end{annexitem}
